\section{Threat model}
\label{sec:threat}

\sysname runs above the Linux kernel as a userspace runtime shared by
mutually distrusting components on a single host. We draw the trust
boundary explicitly here, because everything the paper claims depends
on where it sits.

\para{Trusted}
\begin{itemize}
\itemsep0pt
\item The Linux kernel and its memory-protection guarantees. We rely
on \texttt{mmap(MAP\_SHARED)} backed by \texttt{memfd\_create}, on
\texttt{seccomp-BPF} for syscall filtering, and on
\texttt{CLONE\_NEWUSER}/\texttt{CLONE\_NEWNS} for namespace isolation.
\item The runtime's own pages: the \texttt{CapabilityManager} slot
pool, the \texttt{ChannelControlBlock} of every established channel,
and the M-21 constant-time NASM primitives.
\item RDRAND as a uniform randomness source. Each token consumes 96
bits of it; a compromised RDRAND shrinks the forgery space
accordingly.
\end{itemize}

\para{Untrusted}
\begin{itemize}
\itemsep0pt
\item Any process holding a \captoken. The token's bits may be
anything---we assume the holder can forge, mutate, replay, or simply
guess identifiers.
\item Every byte that crosses a channel. Parsing untrusted input
(framing, headers, sequence numbers) must never push the runtime into
a state-altering action.
\item The peer endpoint. A capability-bearing producer may misbehave:
publish garbage, withhold messages, flood. The consumer's gate
protects the consumer from the producer; the producer's gate protects
\sysname from a producer trying to publish past its own revocation.
\end{itemize}

\para{In scope: forgery, replay, use-after-revoke}
Validation must reject
\textit{(a)} a token whose 128-bit identifier is not currently active
in the pool (forgery);
\textit{(b)} a token whose identifier was once active but has been
revoked, including through cascading revocation of an ancestor
(use-after-revoke); and
\textit{(c)} a token with a correct identifier whose permission mask
does not cover the requested operation (privilege escalation).
The slot-indexed validate path (\Cref{sec:design:fastpath}) does not
weaken any of this: the slot index in the high 32 bits is a hint, not
a credential, and the constant-time compare over the full 128 bits
still has to match.

\para{In scope: side channels in the gate}
The comparison must not reveal which byte differed first, and the
runtime must not take observably different paths based on token
contents beyond the accept/reject bit. The 128-bit identifier match
uses the M-21 \texttt{asm\_ct\_compare} primitive, which is
constant-time and barrier-balanced. The slot index \emph{is} read
before the constant-time compare, but only to pick a pool offset; an
attacker steering the slot index cannot induce distinguishable cache
behaviour across pool sizes (\Cref{sec:eval:perfctrs}).

\para{Out of scope: rowhammer-class faults, kernel TCB bypass}
We do not defend against rowhammer corruption of the token pool,
against speculative execution that pierces the kernel TCB (HongMeng's
territory~\cite{hongmeng}), or against an attacker who can already
write \sysname's process memory. If Linux process isolation fails,
so do we.

\para{Out of scope: cross-host delegation}
This paper is single-host. Cross-host delegation, attestation
handshakes, and remote token transfer belong to follow-up work
(M-23 ConfCompute, in preparation~\cite{astra-confcompute}).

\para{Out of scope: hardware capability enforcement}
We assume no CHERI, Morello, MTE, or LAM. CAP-VMs~\cite{capvms} lives
at the same conceptual layer but needs CHERI silicon; making that
layer available on commodity x86-64 and aarch64 is precisely the
point of \sysname.
